\subsection{TENSOR PRODUCT OF TWO MODULES}

Let \(G_{1}, G_{2}\) be two Z-modules; a mapping u of the set \(G = G_{1} \times G_{2}\) into a Z-module is called biadditive (or Z-bilinear) if \(u(x_{1}, x_{2})\) is ``additive with respect to \(x_{1}\) and with respect to \(x_{2}\) ''; to be precise, this means that, for \(x_{1}, y_{1}\) in \(G_{1}, x_{2}, y_{2}\) in \(G_{2}\) ,

\[
\begin{array}{l} u (x _ {1} + y _ {1}, x _ {2}) = u (x _ {1}, x _ {2}) + u (y _ {1}, x _ {2}) \\ u (x _ {1}, x _ {2} + y _ {2}) = u (x _ {1}, x _ {2}) + u (x _ {1}, y _ {2}). \end{array}
\]

Note that this implies in particular that \(u(0, x_2) = u(x_1, 0) = 0\) for all \(x_1 \in \mathbf{G}_1, x_2 \in \mathbf{G}_2\) .

Let A be a ring, E a right A-module and F a left A-module. We are going to consider the universal mapping problem (Set Theory, IV, § 3, no. 1) where \(\Sigma\) is the species of Z-module structure (the morphisms then being Z-linear mappings, in other words, additive group homomorphisms) and the \(\alpha\) -mappings are the mappings \(f\) of E \(\times\) F into a Z-module G which are Z-bilinear and further satisfy, for all \(x \in \mathbf{E}\) , \(y \in \mathbf{F}\) and \(\lambda \in \mathbf{A}\)

\[
f (x \lambda , y) = f (x, \lambda y).\tag{1}
\]

We show that this problem admits a solution. For this we consider the Z-module \(C = \mathbf{Z}^{(E \times F)}\) of formal linear combinations of the elements of \(E \times F\) with coefficients in Z (§ 1, no. 11), a basis of which can be considered to consist of the ordered pairs \((x, y)\) , where \(x \in E\) and \(y \in F\) . Let D be the sub-Z-module of C generated by the elements of one of the following types:

\[
\left\{ \begin{array}{l} (x _ {1} + x _ {2}, y) - (x _ {1}, y) - (x _ {2}, y) \\ (x, y _ {1} + y _ {2}) - (x, y _ {1}) - (x, y _ {2}) \\ (x \lambda , y) - (x, \lambda y) \end{array} \right.\tag{2}
\]

where \(x, x_{1}, x_{2}\) are in E, \(y, y_{1}, y_{2}\) are in F and \(\lambda\) is in A.

DEFINITION 1. The tensor product of the right A-module E and the left A-module F, denoted by \(\mathbf{E} \otimes_{\mathbf{A}} \mathbf{F}\) or \(\mathbf{E} \otimes_{\mathbf{A}} \mathbf{F}\) (or simply \(\mathbf{E} \otimes \mathbf{F}\) if no confusion is to be feared) is the quotient \(\mathbf{Z}\) -module C/D (the quotient of the \(\mathbf{Z}\) -module C of formal linear combinations of elements of \(\mathbf{E} \times \mathbf{F}\) with coefficients in \(\mathbf{Z}\) , by the submodule \(\mathbf{D}\) generated by the elements of one of the types (2)). For \(x \in \mathbf{E}\) and \(y \in \mathbf{F}\) , the element of \(\mathbf{E} \otimes_{\mathbf{A}} \mathbf{F}\) which is the canonical image of the element \((x, y)\) of \(\mathbf{C} = \mathbf{Z}^{(\mathbf{E} \times \mathbf{F})}\) is denoted by \(x \otimes y\) and called the tensor product of \(x\) and \(y\) .

The mapping \((x,y)\mapsto x\otimes y\) of \(\mathbf{E}\times \mathbf{F}\) into \(\mathbf{E}\otimes_{\mathbf{A}}\mathbf{F}\) is called canonical. It is a Z-bilinear mapping which satisfies conditions (1).

We show that the tensor product \(E \otimes_{A} F\) and the above canonical mapping form a solution of the universal mapping problem posed earlier. To be precise:

PROPOSITION 1. (a) Let \(g\) be a \(\mathbf{Z}\) -linear mapping of \(\mathbf{E} \otimes_{\mathbf{A}} \mathbf{F}\) into a \(\mathbf{Z}\) -module \(\mathbf{G}\) . The mapping \((x, y) \mapsto f(x, y) = g(x \otimes y)\) of \(\mathbf{E} \times \mathbf{F}\) into \(\mathbf{G}\) is \(\mathbf{Z}\) -bilinear and satisfies conditions (1).

\begin{enumerate}
\def\labelenumi{(\alph{enumi})}
\setcounter{enumi}{1}
\tightlist
\item
  Conversely, let \(f\) be a \(\mathbf{Z}\) -bilinear mapping of \(\mathbf{E} \times \mathbf{F}\) into a \(\mathbf{Z}\) -module \(\mathbf{G}\) satisfying conditions (1). Then there exists one and only one \(\mathbf{Z}\) -linear mapping \(g\) of \(\mathbf{E} \otimes_{\mathbf{A}} \mathbf{F}\) into \(\mathbf{G}\) such that \(f(x,y) = g(x \otimes y)\) for \(x \in \mathbf{E}, y \in \mathbf{F}\) .
\end{enumerate}

If \(\phi\) denotes the canonical mapping of \(\mathbf{E} \times \mathbf{F}\) into \(\mathbf{E} \otimes_{\mathbf{A}} \mathbf{F}\) , then \(f = g \circ \phi\) ; whence (a). To show (b), we note that, in the notation of Definition 1, \(f\) extends to a \(\mathbf{Z}\) -linear mapping \(\bar{f}\) of \(\mathbf{C}\) into \(\mathbf{G}\) ( \(\S 1\) , no. 11, Proposition 17). By virtue of relations (1), \(\bar{f}\) is zero for all the elements of \(\mathbf{C}\) of one of the types (2) and hence on \(\mathbf{D}\) . There therefore exists a \(\mathbf{Z}\) -linear mapping \(g\) of \(\mathbf{C} / \mathbf{D} = \mathbf{E} \otimes_{\mathbf{A}} \mathbf{F}\) into \(\mathbf{G}\) such that \(\bar{f} = g \circ \psi\) , where \(\psi: \mathbf{C} \to \mathbf{C} / \mathbf{D}\) is the canonical homomorphism ( \(\S 1\) , no. 8, Remark). The uniqueness of \(g\) is immediate since \(\mathbf{E} \otimes_{\mathbf{A}} \mathbf{F}\) is generated, as a \(\mathbf{Z}\) -module, by the elements of the form \(x \otimes y\) .

Proposition 1 defines a canonical isomorphism of the Z-module of Z-bilinear mappings f of E × F into G, satisfying conditions (1), onto the Z-module Hom \(_{Z}\) (E ⊗A F, G).

When A = Z, conditions (1) are automatically satisfied for every Z-bilinear mapping f and the submodule D of C is already generated by the elements of the first two types in (2).

If now we return to the general case and \(E'\) and \(F'\) denote the underlying Z-modules of E and F respectively, the above remark and Definition 1 show immediately that the Z-module \(E \otimes_{A} F\) can be canonically identified with the quotient of the Z-module \(E' \otimes_{Z} F'\) by the sub-Z-module generated by the elements of the form \((x\lambda) \otimes y - x \otimes (\lambda y)\) , where x runs through E, y runs through F and \(\lambda\) runs through A.

COROLLARY 1. Let H be a Z-module and \(h:E \times F \rightarrow H\) a Z-bilinear mapping satisfying conditions (1) and such that H is generated by \(h(E \times F)\) . Suppose that for every Z-module G and every Z-bilinear mapping f of \(E \times F\) into G satisfying (1)

there exists a Z-linear mapping \(g:H \to G\) such that \(f = g \circ h\) . Then, if \(\phi\) denotes the canonical mapping of \(E \times F\) into \(E \otimes_{A} F\) , there exists one and only one isomorphism \(\theta\) of \(E \otimes_{A} F\) onto H such that \(h = \theta \circ \phi\) .

The hypothesis that \(h(\mathbf{E} \times \mathbf{F})\) generates H implies the uniqueness of g; the corollary is then just the general uniqueness property of a solution of a universal mapping problem (Set Theory, IV, § 3, no. 1).

COROLLARY 2. Let \(\mathbf{E}^0\) (resp. \(\mathbf{F}^0\) ) denote the module \(\mathbf{E}\) (resp. \(\mathbf{F}\) ) considered as a left (resp. right) module over the opposite ring \(\mathbf{A}^0\) ; then there exists one and only one \(\mathbf{Z}\) -module isomorphism \(\sigma: \mathbf{E} \otimes_{\mathbf{A}} \mathbf{F} \to \mathbf{F}^0 \otimes_{\mathbf{A}^0} \mathbf{E}^0\) such that \(\sigma(x \otimes y) = y \otimes x\) for \(x \in \mathbf{E}\) and \(y \in \mathbf{F}\) (``commutativity'' of tensor products).

By definition of the \(A^{0}\) -module structures on \(E^{0}\) and \(F^{0}\) , the mapping \((x,y)\mapsto y\otimes x\) of \(E\times F\) into \(F^{0}\otimes_{A^{0}}E^{0}\) is Z-bilinear and satisfies conditions (1), whence the existence and uniqueness of the Z-linear mapping \(\sigma\) . Similarly a Z-linear mapping \(\tau:F^{0}\otimes_{A^{0}}E^{0}\to E\otimes_{A}F\) is defined such that \(\tau(y\otimes x)=x\otimes y\) and clearly \(\sigma\) and \(\tau\) are inverse isomorphisms.

Remark. The tensor product of non-zero modules may be zero: for example, taking the two Z-modules E = Z/2Z and F = Z/3Z, 2x = 0 and 3y = 0 for all \(x \in E\) and \(y \in F\) ; therefore, in E \(\otimes_{\mathbf{Z}} F\) ,

\[
x \otimes y = 3 (x \otimes y) - 2 (x \otimes y) = x \otimes (3 y) - (2 x) \otimes y = 0
\]

for all \(x\) and \(y\) (cf.~no. 6, Corollary 4 to Proposition 6).

